% CV Simon Zweers -- compile with: pdflatex cv.tex  (or: latexmk -pdf cv.tex)
\documentclass[10pt]{article}

% ---------- Page layout ----------
\usepackage[letterpaper,top=0.55in,bottom=0.6in,left=0.6in,right=0.6in]{geometry}
\usepackage{paracol}
\setcolumnwidth{0.58\textwidth,0.34\textwidth}
\setlength{\columnsep}{0.08\textwidth}
\pagestyle{empty}
\setlength{\parindent}{0pt}
\setlength{\parskip}{0pt}

% ---------- Fonts (pdfLaTeX-friendly stand-ins for Merriweather / Open Sans) ----------
\usepackage[T1]{fontenc}
\usepackage[utf8]{inputenc}
\usepackage{charter}            % serif body, close to Merriweather
\usepackage[sfdefault=false]{roboto}       % sans for small labels, close to Open Sans
\renewcommand{\familydefault}{\rmdefault}
\usepackage{microtype}
\usepackage{fontawesome5}

% ---------- Colors ----------
\usepackage{xcolor}
\definecolor{accent}{HTML}{2079C7}
\definecolor{bodytext}{HTML}{555555}
\definecolor{datetext}{HTML}{777777}
\definecolor{heading}{HTML}{000000}

% ---------- Links ----------
\usepackage[hidelinks]{hyperref}
\hypersetup{pdftitle={CV Simon Zweers},pdfauthor={Simon Zweers}}

\usepackage{tikz}

% ---------- Commands ----------
\newcommand{\cvname}[1]{{\fontsize{30}{34}\selectfont\bfseries\color{heading}#1\par}\vspace{6pt}}

\newcommand{\cvsection}[1]{%
  \vspace{18pt}%
  {\sffamily\bfseries\footnotesize\color{accent}\MakeUppercase{#1}\par}%
  \vspace{12pt}}

% \cventry{Organisation}{Place}{Role}{Dates}{Description}
\newcommand{\cventry}[5]{%
  {\raggedright\large\color{heading}\textbf{#1,} #2 --- \textit{#3}\par}%
  \vspace{3pt}%
  {\sffamily\scriptsize\color{datetext}\MakeUppercase{#4}\par}%
  \vspace{3pt}%
  {\raggedright\small\color{bodytext}#5\par}%
  \vspace{12pt}}

% \skillgroup{Category}{Comma-separated skills}
\newcommand{\skillgroup}[2]{%
  {\raggedright\small\bfseries\color{heading}#1\par}%
  \vspace{2pt}%
  {\raggedright\small\color{bodytext}#2\par}%
  \vspace{10pt}}

\newcommand{\skill}[1]{{\raggedright\small\color{bodytext}#1\par}\vspace{7pt}}

% Contact line: icon + clickable, visibly-styled link
\newcommand{\contact}[3]{%
  \makebox[1.4em][l]{\color{accent}#1}\href{#2}{\color{accent}\underline{#3}}\par\vspace{3pt}}

% Coloured placeholder the size of the photo
\newcommand{\photobox}{%
  \begin{tikzpicture}
    \fill[accent!35, rounded corners=4pt] (0,0) rectangle (1.65in,1.58in);
  \end{tikzpicture}\par}

% =====================================================================
\begin{document}
\begin{paracol}{2}

% ---------------- Left column: header ----------------
\vspace*{4pt}
\cvname{Simon Zweers}
{\raggedright\small\color{bodytext}
Ik ben een enthousiaste en veelzijdige student Electrical Engineering aan de
Technische Universiteit Eindhoven, met een passie voor zowel muziek als
technologie. Al van jongs af aan speel ik cello en heb ik ervaring als
orkestmusicus. Deze muzikale achtergrond heeft mijn discipline, teamwork
en creatief denkvermogen veel geholpen. Daarnaast heb ik een langdurige
interesse in computers en technologie, wat me heeft geleid tot mijn huidige
studie, waarin ik mijn analytische vaardigheden en technische kennis
verder ontwikkel.\par}

% ---------------- Right column: photo + contact ----------------
\switchcolumn
\photobox
\vspace{8pt}
{\sffamily\footnotesize
\contact{\faPhone*}{tel:+31683426852}{+31 6 83426852}
\contact{\faEnvelope}{mailto:simon.zweers@protonmail.com}{simon.zweers@protonmail.com}
\contact{\faLinkedin}{https://www.linkedin.com/in/simon-zweers-754ba0284}{in/simon-zweers}
\contact{\faGlobe}{https://simonzweers.com}{simonzweers.com}
}

% ---------------- Left: experience ----------------
\switchcolumn*
\cvsection{Ervaring}

\cventry{Hogeschool van Amsterdam}{Amsterdam}{Student-onderzoeker}{Juli 2024 - Februari 2026}
  {Als student-onderzoeker ben ik verder gegaan aan mijn stageproject (zie hieronder voor meer informatie).}

\cventry{SURF}{Amsterdam}{Stagiair}{Februari 2024 - Juli 2024}
  {Bij SURF heb ik in het tweede jaar van mijn studie Technische Informatica
   stage gelopen. Ik heb bij SURF gewerkt aan het ontwikkelen van een IoT
   edge applicatie voor het verzamelen van sensordata, in samenwerking met
   de Hogeschool van Amsterdam.}

\cventry{Studentaanhuis}{Amsterdam}{Student ondersteuner}{Augustus 2023 - Februari 2024}
  {Bij Studentaanhuis heb ik vooral technische ondersteuning geboden aan
   ouderen. Dit zijn mensen die vaak extra hulp nodig hadden bij het gebruik
   of begrijpen van een stuk software, hun computer of andere technische
   problemen.}

% ---------------- Right: skills + languages ----------------
\switchcolumn
\vspace{24pt}
\cvsection{Vaardigheden}

\skillgroup{(Embedded Systems) Software}
  {C, C++, Make, CMake, Java, Arduino, ESP32, PlatformIO, Python}
\skillgroup{Hardware \& Prototyping}
  {KiCAD, LTSpice, Qucs-s}
\skillgroup{Web/Internet of Things}
  {MQTT, HTTP, HTML/CSS/JavaScript, Go, Python, MongoDB, MySQL, Docker}
\skillgroup{Tools \& Others}
  {Git, GitHub, GitLab, Figma, Linux, FreeCAD, LaTeX, MATLAB, Simulink}

% \cvsection{Talen}
% \skill{Nederlands}
% \vspace{6pt}
% \skill{Engels}

\end{paracol}

% =====================================================================
\clearpage
\begin{paracol}{2}
\cvsection{Opleiding}

\cventry{Technische Universiteit Eindhoven}{Eindhoven}{Electrical Engineering}{2025 - Heden}
  {In 2025 ben ik begonnen met de universitaire opleiding ``Electrical
   Engineering'' aan de Technische Universiteit Eindhoven.}

\cventry{Hogeschool van Amsterdam}{Amsterdam}{Technische Informatica}{2023 - 2025}
  {Na mijn propedeuse gehaald te hebben, ben ik verder gegaan met mijn
   bachelor Technische Informatica.}

\cventry{Hogeschool van Amsterdam}{Amsterdam}{Propedeuse Technische Informatica (Cum Laude)}{2022 - 2023}
  {Ik heb in 2023 mijn propedeuse voor de opleiding HBO-ICT: Technische
   Informatica met het predicaat ``cum laude'' behaald.}


\switchcolumn
\end{paracol}
\end{document}
